% Generated draft: requires mathematical and source review.
\documentclass[11pt]{article}
\usepackage[margin=25mm]{geometry}
\usepackage{fontspec}
\setmainfont{Latin Modern Roman}
\usepackage{amsmath,amssymb,mathtools,unicode-math}
\setmathfont{Latin Modern Math}
\usepackage{enumitem,needspace,xurl,hyperref}
\setlength{\parindent}{0pt}
\setlength{\parskip}{5pt}
\setlength{\emergencystretch}{3em}
\newcommand{\N}{\mathbb{N}}
\newcommand{\Z}{\mathbb{Z}}
\newcommand{\Q}{\mathbb{Q}}
\newcommand{\R}{\mathbb{R}}
\newcommand{\C}{\mathbb{C}}
\newcommand{\F}{\mathbb{F}}
\newcommand{\A}{\mathbb{A}}
\newcommand{\PP}{\mathbb{P}}
\newcommand{\E}{\mathbb{E}}
\newcommand{\eps}{\varepsilon}
\newcommand{\dd}{\,\mathrm{d}}
\begin{document}
\section*{UnsolvedMath ID: 30001941}
\textbf{Source problem code:} catalogue entry
\subsection{Definitions and mathematical statement}
Let \(u(t,\cdot)\) denote the fast-diffusion solution referred to in the supplied record, and let \(\mathcal F_\sigma\) denote its free-energy functional. Write, wherever this functional is defined,
\[
\mathcal E(t):=\mathcal F_\sigma[u(t,\cdot)],\qquad t\in[0,\infty).
\]
The source question is to obtain sharp estimates for the decay of \(\mathcal E(t)\) valid for every \(t\geq0\), including \(t=0\).

The supplied record does not specify the fast-diffusion equation or its normalization, the spatial domain or dimension, the diffusion exponent, initial or boundary conditions, the admissible solution class, the definition of \(\mathcal F_\sigma\), or the meaning and possible time dependence of \(\sigma\). It also does not specify what constitutes a sharp estimate, or which parameters and initial-data quantities may enter it. Consequently, this is a source-level question with unresolved hypotheses and definitions; it does not determine a precise quantitative inequality or a universally quantified conjecture over a specified class of solutions.
\subsection{Short English statement}
Determine sharp bounds for the decay of the free energy \(\mathcal F_\sigma[u(t,\cdot)]\) of the intended fast-diffusion solution, valid at every time \(t\geq0\). The equation, admissible data, functional, parameter \(\sigma\), and criterion of sharpness require clarification from the source before a precise mathematical formulation is possible.
\subsection{Sources}
\begin{itemize}
\item Local UnsolvedMath registry, permanent ID \(30001941\), title: \emph{Sharp Free-Energy Decay for Fast Diffusion}; source problem code \texttt{catalogue entry}; collection: \emph{Oberwolfach Reports Open Problems}. The supplied record lists proposed year 2011 and no named proposer.
\item Supplied report citation: \emph{Variational Methods for Evolution} (2012), Oberwolfach Reports, DOI 10.4171/owr/2011/55, pp. 3145--3216. \url{https://doi.org/10.4171/owr/2011/55}. The background describes Oberwolfach Reports as records of mathematical workshops and their open-problem discussions. No problem-specific page or passage is supplied.
\item Supplied literature source, labeled ``original OWR report'': \url{https://doi.org/10.4171/owr/2011/55}.
\item Literature assessment dated exactly \texttt{2026-08-21}: the supplied curation labels the problem open and states that the source explicitly calls sharp all-time decay estimates an open question, while later results provide bounds without settling the stated sharp problem. This is supplied dated curation, not independently verified here; no references for those later results are supplied. The background explicitly requests independent verification of time-sensitive status and cited claims.
\item The registry assigns the statement status \texttt{corrected\_verified} and reports that its mathematical content was checked against the official Oberwolfach report, with notation, context, and question form normalized without changing the source-level claim. This verification assertion is supplied provenance and has not been independently checked here.
\end{itemize}
\end{document}
